\section{A classical-polylogarithm generator for every pair}
\label{tail:generator-section}

The finite formula of Theorem~\ref{tail:all-moments} can be summed
at every moment order. This yields an analytic identity, as opposed
to a merely formal generating series.

For integers \(a\ge2\) and \(b\ge1\), define
\begin{equation}\label{tail:I}
 I_{a,b}(t)=\int_0^t u^a(t-u)^b\coth(u/2)\,\dd u.
\end{equation}
Inside \(|t|<2\pi\), use the straight segment from zero to \(t\).
The integrand has a removable zero-endpoint singularity, and the
integral is unambiguous there. Define the finite polynomial
\begin{equation}\label{tail:P}
 P_{p,q}(t)=\sum_{\substack{1\le d<q\\d\ {
 \rm odd}}}
 \left\{\zeta(p,q-d)+\frac12\zeta(W-d)\right\}
                        \frac{t^{d-1}}{d!}.
\end{equation}
Every double zeta in this polynomial is eliminated by
\eqref{tail:Euler-reduction} or \eqref{tail:Euler-boundary}.

\begin{theorem}[All-pair analytic moment generator]
\label{tail:generator}
For even \(p,q\ge2\), the series
\[
 G_{p,q}(t)=\sum_{r=0}^{\infty}
           D_{p,q}(-2r)\frac{t^{2r}}{(2r)!}
\]
converges in \(|t|<2\pi\). In that disk it is the even holomorphic
function
\begin{equation}\label{tail:generator-formula}
 \boxed{
 G_{p,q}(t)=\frac{t}{2\sinh(t/2)}
             \{F_{p,q}(t)+F_{q,p}(t)\},}
\end{equation}
where all singularities at zero are removable and
\begin{align}\label{tail:F}
 F_{p,q}(t)={}&P_{p,q}(t)\notag\\
 &+\sum_{h=1}^{p/2-1}
   \frac{\zeta(p-2h)}{2t(2h)!(q-1)!}
       I_{2h,q-1}(t)
  -\frac{I_{p,q-1}(t)}{4t p!(q-1)!}.
\end{align}
For real \(0<t<2\pi\), the expression in braces has a finite
form using only rational functions of \(t\), ordinary zeta values,
and \(\operatorname{Li}_k(e^{-t})\) for \(2\le k\le p+q\).
The full generator includes the displayed elementary prefactor
\(t/(2\sinh(t/2))\).
The integral formula fixes its holomorphic continuation at zero.
\end{theorem}

\begin{proof}
The Bernoulli convolution in \eqref{tail:moment-formula} is
equivalent formally to
\begin{equation}\label{tail:block-generator}
 G_{p,q}(t)=\frac{t}{2\sinh(t/2)}
           \sum_{\substack{d\ge1\\d\ {\rm odd}}}
                        C_{p,q}(d)\frac{t^{d-1}}{d!}.
\end{equation}
We show that the last series converges in the stated disk and
equals \(F_{p,q}+F_{q,p}\).

Write
\[
 E(u)=\frac u2\coth(u/2)
     =\sum_{k\ge0}\frac{B_{2k}}{(2k)!}u^{2k},
 \qquad |u|<2\pi.
\]
The terms with \(d<q\) in the ordered block are exactly
\(P_{p,q}\). For each fixed \(1\le h\le p/2-1\), put
\(d=q+2h-1+2k\). Its remaining Bernoulli contribution is
\[
 \frac{\zeta(p-2h)}{(2h)!}
 \sum_{k\ge0}\frac{(2h+2k-1)!}{(q+2h+2k-1)!}
            \frac{B_{2k}}{(2k)!}t^{q+2h+2k-2}.
\]
Euler's beta integral transforms this normally convergent series
into
\[
 \frac{\zeta(p-2h)}{t(2h)!(q-1)!}
       \int_0^t u^{2h-1}(t-u)^{q-1}E(u)\,\dd u
   =\frac{\zeta(p-2h)I_{2h,q-1}(t)}
          {2t(2h)!(q-1)!}.
\]
For the first binomial term in \(R_{p,q}(d)\), set
\(d=W-1+2k\). Its series is
\[
 -\frac1{2p!}\sum_{k\ge0}
    \frac{(p+2k-1)!}{(W+2k-1)!}
         \frac{B_{2k}}{(2k)!}t^{W+2k-2},
\]
which the same beta integral turns into
\(-I_{p,q-1}(t)/(4t p!(q-1)!)\). The second term is the
same expression with \(p,q\) exchanged. This proves
\eqref{tail:F} and \eqref{tail:generator-formula}.

All power-series integrations above are normal on compact subsets
of \(|t|<2\pi\). The prefactor \(t/(2\sinh(t/2))\) is
holomorphic there, including at zero. The resulting function is
even: this follows from its Taylor construction, or by substituting
\(u=tv\) in each integral. Thus
\eqref{tail:block-generator} is the Taylor series of an even
holomorphic function, proving the claimed convergence.

Finally, write \(\coth(u/2)=1+2/(e^u-1)\) and expand the
polynomial \((t-u)^b\). For \(t>0\), elementary integration
gives
\begin{align}\label{tail:finite-Li}
 I_{a,b}(t)={}&\frac{a!b!}{(a+b+1)!}t^{a+b+1}\notag\\
 &+2\sum_{\ell=0}^{b}(-1)^\ell\binom b\ell
                     t^{b-\ell}J_{a+\ell}(t),\notag\\
 J_m(t)={}&m!\left\{\zeta(m+1)
       -\sum_{j=0}^{m}\frac{t^j}{j!}
                  \operatorname{Li}_{m+1-j}(e^{-t})\right\}.
\end{align}
One proof expands \((e^u-1)^{-1}\) into its positive exponential
series and integrates each term. Since \(m\ge2\), the
integrals are absolutely convergent. The coefficient of
\(\operatorname{Li}_1(e^{-t})\) in the first two lines is
proportional to \(\sum_{\ell=0}^{b}(-1)^\ell\binom b\ell=0\).
The largest remaining polylogarithm order is
\(a+b+1\le p+q\). This proves the last assertion.
\end{proof}

\subsection{An explicit \texorpdfstring{$\operatorname{Li}_2$--$\operatorname{Li}_6$}{Li2--Li6} identity}

For \((p,q)=(2,4)\), the general generator becomes
\begin{align}\label{tail:24-integral-generator}
 G_{2,4}(t)=\frac{t}{2\sinh(t/2)}\bigg\{
 &\frac{15}{2}\zeta(5)-3\zeta(2)\zeta(3)
       +\frac{t^2}{4}\zeta(3)
       +\frac{\zeta(2)}{4t}I_{2,1}(t)\notag\\
 &-\frac{I_{4,1}(t)}{96t}
       -\frac{I_{2,3}(t)}{48t}\bigg\}.
\end{align}
Put \(L_k(t)=\operatorname{Li}_k(e^{-t})\). An entirely explicit
version, valid for real \(0<t<2\pi\), is
\begin{equation}\label{tail:24-Li-generator}
 G_{2,4}(t)=\frac{t}{2\sinh(t/2)}\,\mathcal L_{2,4}(t),
\end{equation}
where
\begin{align}\label{tail:24-Li-body}
 \mathcal L_{2,4}(t)={}&4\zeta(5)-2\zeta(2)\zeta(3)
       +\frac{t^2}{6}\zeta(3)+\frac{3t}{4}\zeta(4)
       +\frac{t^3}{48}\zeta(2)-\frac{t^5}{1440}\notag\\
 &+\left(\frac{t}{2}\zeta(2)-\frac{t^3}{48}\right)L_2(t)
       +\left(2\zeta(2)-\frac{t^2}{6}\right)L_3(t)\notag\\
 &-tL_4(t)-4L_5(t)
       +\frac{3\zeta(2)}{t}\{L_4(t)-\zeta(4)\}
       +\frac{15}{2t}\{\zeta(6)-L_6(t)\}.
\end{align}
Although individual polylogarithms have logarithmic expansions at
zero, their specified combination is even and holomorphic there.
Equation~\eqref{tail:24-integral-generator}, rather than a separate
choice of logarithm for each term, provides the continuation. The
coefficient of \(t^{2r}/(2r)!\) is exactly the all-order moment in
\eqref{tail:moment-formula} and \eqref{tail:24-blocks}.

\subsection{Scope and next exact questions}

The two-tail problem is closed by a finite Bernoulli formula,
ordinary-zeta elimination, and a convergent classical-polylogarithm
generator. The next questions are structurally different.

\begin{enumerate}
\item For four even tails, determine the smallest useful
multiple-zeta coordinate space uniform in the centered moment
order. Finite summation by parts still lowers the number of tails,
but higher-depth constants can enter. A two-tail calculation does
not prove their elimination.
\item Determine exact first spectral derivatives at the regular
points \(s=-2r\). These contain logarithmically weighted remainder
sums in addition to local Bernoulli data. The value theorem alone
does not reduce those global contributions.
\item Introduce rational inner shifts while retaining a matched
center for the outer weight. Identify the resulting finite
colored Euler-sum coordinates and the precise analogues of the
classical-polylogarithm generator.
\item For fixed \(p,q\), find optimal finite generating sets of
the actual moment sequence. Corollary~\ref{tail:fixed-space}
provides at most \(\max\{p,q\}\) explicit coordinates, and
\((2,4)\) has an exact four-moment generator. Arithmetic
independence of a proposed minimal set must be distinguished from
the rank of its rational coefficient matrix.
\end{enumerate}

